% ===========================================================================
\subsection{\file{.github/workflows-disabled/deploy.yml.disabled}}
% ===========================================================================

\textbf{YAML, 103 lines.} The automation that used to publish this site: it
watched for pushes, checked for a marker in the commit message, copied the files
to the agency's server over an encrypted connection, and then checked that the
site still answered. It is currently switched off.

\begin{provenance}{.github/workflows-disabled/deploy.yml.disabled}
Hand-written, added in commit \code{d49e034} on 2026-07-28 at the path
\file{.github/workflows/deploy.yml}, rewritten in \code{2c66a7a} the same day when
the deploy target changed from Netlify to the VPS, fixed three times on 2026-07-29,
and finally renamed to its present path in \code{da508a7}.

Its shape is shared with the sibling project \code{salmanadnan-com}, whose
workflow was written in the same week and carries the same retry loops and the
same \code{[deploy]} marker check. Neither file says which came first, and this
document does not guess. The \code{[deploy]} convention itself is
workspace-wide and predates both. \textbf{Confidence: certain for hand-written and
for the rename history, which is recorded in git; the relationship to the sibling
repository is stated as a resemblance and not as a lineage.}
\end{provenance}

\subsubsection*{Why it exists}

To remove a manual step. Without it, publishing means somebody connecting to the
server and copying files by hand, which is slow and easy to get wrong. With it,
publishing is a consequence of writing a commit message in a particular way.

\begin{jargon}{YAML}
A file format for configuration, where structure is expressed by indentation
rather than by brackets. Two spaces of indent mean ``this belongs to the line
above''. It is unforgiving about indentation, which is its main drawback: a single
misplaced space changes the meaning.
\end{jargon}

\begin{jargon}{GitHub Actions, and a workflow}
A service that runs jobs on GitHub's computers when something happens in a
repository. The jobs are described in a workflow file in
\file{.github/workflows/}. Each job runs on a fresh, empty machine, which is why
the file below has to install its own SSH key every time.
\end{jargon}

\subsubsection*{Line by line}

\paragraph{Lines 1 to 19: the header comment and the trigger.}

\begin{lstlisting}[firstnumber=1,caption={\file{deploy.yml.disabled}, lines 1--19}]
# ai.digitalise.agency: deploy to VPS when commit subject starts with [deploy].
#
# A push to main always checks the marker. It DEPLOYS only when the commit
# subject line STARTS with [deploy], matching every other agency repository.
# Mentioning the marker anywhere else in the message does nothing.
#
# The site is served by Caddy from /root/ai.digitalise.agency on the VPS.
# Caddy picks up the new files immediately, no restart needed.
#
# Required repository secrets:
#   VPS_HOST     46.202.194.85
#   VPS_USER     deploy
#   VPS_SSH_KEY  private key for that user
name: Deploy ai.digitalise.agency

on:
  push:
    branches: [main]
\end{lstlisting}

Lines 1 to 13 are a comment block, and it is the best documentation in the
repository. In YAML a \code{\#} starts a comment. It states the marker rule twice
and in capitals, which tells you it had caught somebody out; it names the serving
path and the fact that Caddy needs no restart; and it lists the three secrets the
workflow needs.

\begin{honest}{The server's address is written in the comment that explains it is
a secret}
Line 11 records \code{VPS\_HOST 46.202.194.85} in plain text, immediately under the
heading ``Required repository secrets''. Whatever the value of keeping the address
in GitHub's secret store, writing it here defeats it.

Proportion matters: an IP address is not a credential, this repository is private,
and the server is reachable on the public internet anyway. The genuine cost is that
it is now in the git history permanently, so if the server is ever replaced this
line becomes stale and misleading rather than merely redundant. The private key,
which is the thing that actually matters, is correctly absent.
\end{honest}

Line 14 names the workflow, which is the label shown in GitHub's interface.

Lines 16 to 18 set the trigger: run on every push to the branch \code{main}, and
nothing else. Not on pull requests, not on a schedule, not manually. That last
absence is worth noting: there is no \code{workflow\_dispatch}, so there is no way
to publish without making a commit.

\paragraph{Lines 20 to 38: the marker check.}

\begin{lstlisting}[firstnumber=19,caption={\file{deploy.yml.disabled}, lines 19--38}]

jobs:
  check-deploy:
    name: Check Deploy Marker
    runs-on: ubuntu-latest
    outputs:
      should_deploy: ${{ steps.check.outputs.should_deploy }}
    steps:
      - id: check
        env:
          COMMIT_MSG: ${{ github.event.head_commit.message }}
        run: |
          SUBJECT="${COMMIT_MSG%%$'\n'*}"
          if [[ "$SUBJECT" == "[deploy]"* ]]; then
            echo "should_deploy=true" >> "$GITHUB_OUTPUT"
            echo "Deploy marker found on the subject line."
          else
            echo "should_deploy=false" >> "$GITHUB_OUTPUT"
            echo "No deploy marker: skipping deploy."
          fi
\end{lstlisting}

The first job exists only to answer a yes-or-no question, and it runs on a fresh
Ubuntu machine to do it. Lines 24 and 25 declare that this job produces an output
called \code{should\_deploy}, which the second job will read.

Line 28 passes the commit message into the script as an environment variable rather
than pasting it directly into the script text. That distinction is a genuine
security measure.

\begin{keyidea}{Why the commit message goes through an environment variable}
If line 31 had been written as \code{SUBJECT="\$\{\{ github.event.head\_commit.message
\}\}"}, GitHub would paste the commit message directly into the script before the
shell ever sees it. Anyone who could write a commit message could then write shell
commands into it and have them run on the machine holding the deploy key.

Passing it as an environment variable and reading it as \code{\$COMMIT\_MSG} means
the message is data, never code. This is a well-known class of vulnerability in
GitHub Actions, and this file gets it right.
\end{keyidea}

Line 31 extracts the subject line: \code{\$\{COMMIT\_MSG\%\%\$'\textbackslash n'*\}}
is shell syntax meaning ``delete the longest match of a newline followed by anything,
from the end'', which leaves everything before the first newline.

Line 32 is the test itself. \code{"\$SUBJECT" == "[deploy]"*} compares with a
trailing wildcard, which is why the marker must \emph{begin} the subject.

Lines 33 and 36 write the answer to \code{\$GITHUB\_OUTPUT}, a file GitHub reads to
collect a step's outputs. Lines 34 and 37 print a human-readable explanation into
the log, which is a small kindness to whoever is wondering why nothing deployed.

\paragraph{Lines 40 to 63: starting the deploy, and the SSH setup.}

\begin{lstlisting}[firstnumber=39,caption={\file{deploy.yml.disabled}, lines 39--63}]

  deploy:
    name: Deploy to VPS
    runs-on: ubuntu-latest
    needs: check-deploy
    if: needs.check-deploy.outputs.should_deploy == 'true'
    steps:
      - uses: actions/checkout@v4

      - name: Authorize SSH
        env:
          KEY: ${{ secrets.VPS_SSH_KEY }}
          HOST: ${{ secrets.VPS_HOST }}
        run: |
          mkdir -p ~/.ssh && chmod 700 ~/.ssh
          printf '%s\n' "$KEY" > ~/.ssh/deploy_key
          chmod 600 ~/.ssh/deploy_key
          for i in 1 2 3 4 5; do
            if ssh-keyscan -T 10 -H "$HOST" >> ~/.ssh/known_hosts 2>/tmp/keyscan.err; then
              [ -s ~/.ssh/known_hosts ] && break
            fi
            echo "ssh-keyscan attempt $i failed: $(cat /tmp/keyscan.err 2>/dev/null)"
            sleep 5
          done
          [ -s ~/.ssh/known_hosts ] || { echo "::error::could not fetch the host key"; exit 1; }
\end{lstlisting}

Line 43 makes this job wait for the first, and line 44 runs it only if the answer
was yes. Everything below is therefore skipped entirely on an ordinary commit.

Line 46 checks out the repository onto the machine, using GitHub's own action.
Without this the machine would be empty.

Lines 53 to 55 install the private key: create the \file{.ssh} folder,
write the key into it, and set permissions so only the owner can read it.
\code{chmod 600} is not optional; SSH refuses to use a key file that other users
can read, and its error message when it does is famously unhelpful.

\begin{jargon}{SSH key pair, and the host key}
Connecting securely involves two separate keys and confusing them is common. The
\emph{deploy key} here is the private half of a pair proving \emph{we} are allowed
in; its public half sits on the server. The \emph{host key} is the server's own
identity, which the client checks to be sure it is talking to the right machine and
not an impostor. Lines 53 to 55 install the first. Lines 56 to 63 fetch the second.
\end{jargon}

Lines 56 to 62 are the retry loop, and they exist because this step failed
repeatedly. \code{ssh-keyscan} asks the server for its host key, with a 10-second
timeout, appending the result to \file{known\_hosts} and any complaint to a
temporary file. \code{[ -s ~/.ssh/known\_hosts ]} then tests that the file is
non-empty, because \code{ssh-keyscan} can succeed as a command while producing
nothing at all. That distinction is the subtle part: checking the exit code alone
would not have been enough.

Line 63 fails the whole job with a message if five attempts produced nothing.
\code{::error::} is GitHub's syntax for marking a line in the log as an error so it
surfaces in the interface.

\begin{keyidea}{Why fetching the host key each time is a real trade-off}
Fetching a server's identity at the moment you connect, and trusting whatever comes
back, means you are not actually verifying anything: an attacker positioned between
the runner and the server could supply their own. The secure alternative is to store
the expected host key as a secret and compare.

That is a real weakness and it is a widespread convention, and the sibling project
\code{passbolt-digitalise} in this workspace does it the stricter way, keeping a
\file{.github/known\_hosts} file in the repository. Two projects in one workspace
solving the same problem two ways is worth knowing about.
\end{keyidea}

\paragraph{Lines 65 to 83: copying the files.}

\begin{lstlisting}[firstnumber=64,caption={\file{deploy.yml.disabled}, lines 64--83}]

      - name: Sync files
        env:
          HOST: ${{ secrets.VPS_HOST }}
          USER: ${{ secrets.VPS_USER }}
        run: |
          retry() {
            for attempt in 1 2 3; do
              "$@" && return 0
              echo "attempt $attempt failed: $*"
              [ "$attempt" -lt 3 ] && sleep $(( attempt * 15 ))
            done
            return 1
          }
          retry rsync -az --no-owner --no-group --omit-dir-times --delete \
            -e "ssh -i ~/.ssh/deploy_key -o BatchMode=yes -o ConnectTimeout=25" \
            --exclude '.git' \
            --exclude '.github' \
            --exclude 'assets/ASSETS.md' \
            ./ "$USER@$HOST:/root/ai.digitalise.agency/"
\end{lstlisting}

Lines 70 to 77 define a small shell function that runs whatever it is given, up to
three times, waiting longer between each attempt: 15 seconds, then 30. \code{"\$@"}
means ``all the arguments, run as a command''. This is the second retry loop in the
file, and like the first it is evidence of a step that failed in practice.

Line 78 is the actual copy. Each flag earns its place:

\begin{table}[H]\centering\small
\begin{tabular}{@{}ll@{}}
\toprule
Flag & What it does \\
\midrule
\code{-a} & archive mode: recurse, and preserve timestamps, permissions and links \\
\code{-z} & compress during transfer \\
\code{--no-owner} & do not try to set file ownership on the far end \\
\code{--no-group} & do not try to set group ownership either \\
\code{--omit-dir-times} & do not try to set directory timestamps \\
\code{--delete} & remove files on the server that no longer exist here \\
\bottomrule
\end{tabular}
\caption{The \code{rsync} flags, and why each is present.}
\end{table}

The three \code{--no}/\code{--omit} flags are all working around the same thing: the
\code{deploy} user is not the owner of the target directory and lacks the privilege
to change ownership or directory timestamps. Without them \code{rsync} reports
errors and fails. These flags are the surviving trace of the problem that commit
\code{f698421} was fixing.

\code{--delete} deserves respect: it makes the server an exact mirror, so removing a
file here removes it there. That is usually what you want and it is also the flag
that turns a mistaken source path into a wiped site.

Line 79 tells \code{rsync} how to connect: use this key, never prompt for a password
(\code{BatchMode=yes}, essential because there is no human to answer), and give up
after 25 seconds.

Lines 80 to 82 exclude three things from the copy: the git history, the workflow
folder, and the internal checklist.

Line 83 is the source and destination. \code{./} means the current directory's
\emph{contents}, and the trailing slash is significant: without it, \code{rsync}
would copy the directory itself into the target, producing a nested folder.

\begin{honest}{This line would now break the live site}
\code{./} was the right source when \file{index.html} sat at the repository root.
Since commit \code{da508a7} moved everything down one level, \code{./} now contains
a single folder called \file{new-design-reference}. Restoring this workflow
unchanged would therefore copy that folder to the server and, because
\code{--delete} is set, remove the real site's files in the same operation.

The correct source after the move would be \code{new-design-reference/}. The same
commit that moved the files disabled this workflow, so the two are consistent; the
trap only exists for somebody who re-enables one without reading the other.
\end{honest}

\begin{honest}{The exclusion for the checklist no longer matches}
\code{--exclude 'assets/ASSETS.md'} was written when that path was correct. The file
now lives at \file{new-design-reference/assets/ASSETS.md}, which this pattern does
not match, so a restored workflow would publish the internal checklist, including its
line about replacing placeholder numbers, to the public internet.
\end{honest}

\paragraph{Lines 85 to 103: verifying, and cleaning up.}

\begin{lstlisting}[firstnumber=84,caption={\file{deploy.yml.disabled}, lines 84--103}]

      - name: Verify
        env:
          HOST: ${{ secrets.VPS_HOST }}
          USER: ${{ secrets.VPS_USER }}
        run: |
          for i in $(seq 1 12); do
            code=$(ssh -i ~/.ssh/deploy_key -o BatchMode=yes "$USER@$HOST" \
              'curl -sf -o /dev/null -w "%{http_code}" --max-time 10 https://ai.digitalise.agency/' || echo 000)
            case "$code" in
              2*|3*) echo "Site answering with $code after ${i} check(s)."; exit 0 ;;
            esac
            echo "waiting, got $code ($i/12)"
            sleep 5
          done
          echo "::warning::site did not respond 2xx/3xx; check Caddy on the VPS"

      - name: Clean up key
        if: always()
        run: rm -f ~/.ssh/deploy_key
\end{lstlisting}

The verification step asks, up to twelve times over about a minute, whether the site
is answering. It does this by connecting to the server and running \code{curl}
\emph{there} rather than from the runner, which tests the server's own view of
itself.

The \code{curl} flags: \code{-s} quiet, \code{-f} treat an error status as a failure,
\code{-o /dev/null} throw the page away, \code{-w "\%\{http\_code\}"} print just the
status number, and a 10-second limit. The \code{|| echo 000} supplies a fallback so
that a failed connection yields the harmless string \code{000} instead of an empty
one.

Line 93 accepts any status beginning with 2 or 3, meaning success or redirect, and
exits successfully.

\begin{honest}{A failed verification is only a warning}
Line 99 uses \code{::warning::} rather than \code{::error::}, and the step does not
exit with a failure code. So if the site never came back, the workflow would still
finish green with a warning buried in the log. Given that the preceding step ran
\code{rsync --delete} against the live site, a failure here is the one moment you
would want a loud red failure and an email. This is the clearest correctness
weakness in the workflow, and it is one word to fix.
\end{honest}

Lines 101 to 103 delete the private key from the runner. \code{if: always()} makes
this run even when an earlier step failed, which is the point: a key must not be left
behind on a machine because something else went wrong. In practice GitHub destroys
the machine anyway, so this is belt and braces, and it is the right instinct.

\subsubsection*{What it connects to}

\textbf{Reads:} three GitHub secrets (\code{VPS\_SSH\_KEY}, \code{VPS\_HOST},
\code{VPS\_USER}) and the commit message of the push that triggered it.

\textbf{Writes:} the contents of \file{/root/ai.digitalise.agency/} on the VPS at
\code{46.202.194.85}, which Caddy serves at \code{https://ai.digitalise.agency/}.

\textbf{Excludes:} \file{.git}, \file{.github} and \file{assets/ASSETS.md}.

\textbf{Runs when:} never, currently. It sits in
\file{.github/workflows-disabled/} with the extension \file{.disabled}, and GitHub
only reads \file{.github/workflows/*.yml}.

\subsubsection*{How it breaks}

\begin{itemize}
\item \textbf{Restored unchanged}, it would copy the wrong directory and, with
  \code{--delete}, take down whatever is currently served.
\item \textbf{A missing or wrong \code{VPS\_SSH\_KEY}} fails at the \code{rsync}
  step after three attempts. This is what commit \code{ed5c154} fixed.
\item \textbf{A \code{deploy} user without permission to traverse into
  \code{/root}} fails the same way. This is what \code{f698421} fixed, and it is the
  least obvious of the three: the user can hold the key, connect successfully, and
  still be unable to reach the target directory.
\item \textbf{GitHub's runner being refused by the server} fails the host-key fetch,
  which is what both retry loops exist for.
\item \textbf{A commit whose subject does not start with \code{[deploy]}} does
  nothing at all, by design, and the log says so.
\item \textbf{The site failing to come back} produces a warning and a green tick,
  which is the failure mode worth fixing first.
\end{itemize}

\subsubsection*{Honest findings for this file}

The server address is recorded in a comment describing it as a secret; the
\code{rsync} source path and the \file{ASSETS.md} exclusion are both stale after
the move; the host key is trusted on first sight rather than pinned, unlike the
sibling \code{passbolt-digitalise} repository; and a failed post-deploy check
produces a warning rather than a failure. Against those: the commit message is
handled safely as data rather than as code, both retry loops are well built, the
key is cleaned up unconditionally, and the comment block at the top is the most
useful documentation in the repository.
